Appendix Table~\ref{tab:preferred-clustering-sensitivity} holds the sample
fixed across three control sets. Its two inference columns use exactly the
same coefficients and different clustering levels. In the preferred model,
sector clustering changes the formal 101+/solo standard error from 0.037
to 0.091 log points, without changing the 16.5 percent estimate.
Table~\ref{tab:preferred-extension-clustering} applies the same comparison
to endpoint changes and elasticities. The informal endpoint change has
\(p=0.011\) under sector-department clustering and \(p=0.090\) under sector
clustering. Both elasticity estimates remain positive and significant
under both conventions.

\begin{table}[htbp]
  \centering
  \caption{Separating controls from clustering in the 101+/solo comparison}
  \label{tab:preferred-clustering-sensitivity}
  \scriptsize
  \begin{tabular}{lrrrrr}
    \toprule
    Controls & Premium (\%) & S.E. (SD) & $p$ (SD) & S.E. (S) & $p$ (S) \\
    \midrule
    \multicolumn{6}{l}{\textit{Formal workers}} \\
    Sector-year FE & 16.6 & 0.038 & $<0.001$ & 0.088 & 0.100 \\
    Sector-department-year FE & 15.8 & 0.039 & $<0.001$ & 0.093 & 0.135 \\
    + Occupational position & 16.5 & 0.037 & $<0.001$ & 0.091 & 0.113 \\
    \midrule
    \multicolumn{6}{l}{\textit{Informal workers}} \\
    Sector-year FE & 39.9 & 0.020 & $<0.001$ & 0.046 & $<0.001$ \\
    Sector-department-year FE & 33.4 & 0.019 & $<0.001$ & 0.043 & $<0.001$ \\
    + Occupational position & 37.1 & 0.019 & $<0.001$ & 0.039 & $<0.001$ \\
    \bottomrule
  \end{tabular}
  \vspace{0.3em}
  \begin{minipage}{0.95\textwidth}
  \footnotesize
  Notes: Specifications from columns (1), (3), and (4) of Table~\ref{tab:firm-size-local-geography-formality-robustness}, estimated on the same 5,268,375 observations. The reference is solo work within each formality group. SD denotes sector-department clustering and S denotes sector clustering. Standard errors are in log points. Two-sided tests use 558 and 16 degrees of freedom, respectively.
  \end{minipage}
\end{table}

\begin{table}[htbp]
  \centering
  \caption{Clustering sensitivity of changes over time and elasticities}
  \label{tab:preferred-extension-clustering}
  \small
  \begin{tabular}{lrrrrr}
    \toprule
    Workers & Estimate & S.E. (SD) & $p$ (SD) & S.E. (S) & $p$ (S) \\
    \midrule
    \multicolumn{6}{l}{\textit{Change in the 101+/solo log-wage gap, 2008--2025}} \\
    Formal & 0.008 & 0.045 & 0.861 & 0.047 & 0.869 \\
    Informal & 0.152 & 0.060 & 0.011 & 0.084 & 0.090 \\
    \midrule
    \multicolumn{6}{l}{\textit{Approximate firm-size elasticity}} \\
    Formal & 0.040 & 0.005 & $<0.001$ & 0.012 & 0.003 \\
    Informal & 0.097 & 0.004 & $<0.001$ & 0.008 & $<0.001$ \\
    \bottomrule
  \end{tabular}
  \vspace{0.3em}
  \begin{minipage}{0.95\textwidth}
  \footnotesize
  Notes: SD denotes sector-department clustering and S denotes sector clustering. Changing clustering leaves coefficients unchanged. The first panel reports the 2025 minus 2008 change in the log-wage contrast from the joint endpoint model, not a percentage-point change. The second panel reports slopes from Equation~(\ref{eq:preferred-elasticity}). Two-sided tests use 539 SD degrees of freedom for the endpoint model, 558 for elasticities, and 16 for sector clustering in both models.
  \end{minipage}
\end{table}

\FloatBarrier

Table~\ref{tab:salaried-robustness} compares the full-sample model
with a regression restricted to salaried workers. Both use 11--19 workers
as the reference within each formality group and retain the same controls,
weights, and fixed-effect structure. The salaried restriction leaves
2,682,075 observations before removing 66 additional singleton
sector-department-year cells, for a final sample of 2,682,009.
Standard errors use 555 sector-department clusters in this restricted
sample, compared with 559 in the full sample. The formal 101+ log contrast
is 0.121 in the restricted sample, equivalent to a 12.9 percent premium
over formal 11--19 workers. The informal contrast is -0.033 log points,
equivalent to a 3.2 percent shortfall relative to informal 11--19 workers.

\begin{table}[!htbp]
  \centering
  \caption{Firm-size premiums in the full and salaried-only samples}
  \label{tab:salaried-robustness}
  \small
  \begin{tabular}{lrrrr}
    \toprule
    & \multicolumn{2}{c}{All workers} & \multicolumn{2}{c}{Salaried workers} \\
    \cmidrule(lr){2-3}\cmidrule(lr){4-5}
    Firm size & Informal & Formal & Informal & Formal \\
    \midrule
    Solo & -0.354*** & -0.041 & -0.226*** & -0.064*** \\
     & (0.018) & (0.031) & (0.016) & (0.023) \\
    2--3 & -0.242*** & -0.116*** & -0.180*** & -0.082*** \\
     & (0.011) & (0.014) & (0.011) & (0.013) \\
    4--5 & -0.098*** & -0.064*** & -0.080*** & -0.058*** \\
     & (0.007) & (0.007) & (0.007) & (0.007) \\
    6--10 & -0.030*** & -0.028*** & -0.028*** & -0.022*** \\
     & (0.006) & (0.007) & (0.005) & (0.006) \\
    11--19 & Reference & Reference & Reference & Reference \\
    20--30 & 0.010 & 0.002 & 0.011* & 0.011*** \\
     & (0.007) & (0.004) & (0.007) & (0.003) \\
    31--50 & 0.009 & 0.023*** & 0.009 & 0.035*** \\
     & (0.010) & (0.005) & (0.010) & (0.004) \\
    51--100 & 0.021 & 0.048*** & 0.030** & 0.062*** \\
     & (0.013) & (0.006) & (0.012) & (0.006) \\
    101+ & -0.040** & 0.118*** & -0.033** & 0.125*** \\
     & (0.017) & (0.014) & (0.017) & (0.012) \\
    \midrule
    11--19/solo premium (\%) & 42.5 & 4.2 & 25.3 & 6.6 \\
    101+/solo premium (\%) & 37.0 & 17.2 & 21.3 & 20.8 \\
    Equal 101+/solo premiums: $p$ & \multicolumn{2}{c}{$<0.001$} & \multicolumn{2}{c}{0.888} \\
    Equal size profiles: $p$ & \multicolumn{2}{c}{$<0.001$} & \multicolumn{2}{c}{$<0.001$} \\
    Observations & \multicolumn{2}{c}{5,426,174} & \multicolumn{2}{c}{2,765,713} \\
    \bottomrule
  \end{tabular}
  \vspace{0.3em}
  \begin{minipage}{0.95\textwidth}
  \footnotesize
  Notes: Main rows report within-formality log-wage contrasts relative to 11--19 workers. Summary premiums use solo work as the reference and are in percent. Each column pair estimates Equation~(\ref{eq:main}) on the indicated sample. The restricted sample excludes 66 additional singleton cells. Standard errors in parentheses use sector-department clustering, with $G=559$ clusters in the full sample and $G=555$ among employees. Equality tests compare formal and informal log contrasts relative to solo work, individually for 101+ (two-sided $t$ test) and jointly across all sizes (Wald $F$ test with eight numerator degrees of freedom). Both use $G-1$ denominator degrees of freedom. * $p<0.10$, ** $p<0.05$, *** $p<0.01$.
  \end{minipage}
\end{table}

\FloatBarrier

The common-premium comparison model is
\begin{equation}\label{eq:common-premium}
  \ln(w_{ist}) = \alpha
  + \sum_{g\neq g_0} b_g D_{ig}
  + X_{ist}'\gamma + \phi F_{ist}
  + \lambda_{s\times t}+\varepsilon_{ist}.
\end{equation}
This model constrains the size coefficients to be the same for formal and
informal workers and omits occupational-position controls.
Table~\ref{tab:firm-size-wage-premium-regressions} starts with size indicators
alone, then adds sector-year effects, worker characteristics, and formality.
Figure~\ref{fig:firm-size-wage-premium} expresses the final column as
percentage premiums. The following alternative-reference tables retain
the same estimates and change only their normalization.

\begin{table}[htbp]
  \centering
  \caption{Common-premium comparison: sequential controls}
  \label{tab:firm-size-wage-premium-regressions}
  \small
  \begin{tabular}{lcccc}
    \toprule
    & \multicolumn{4}{c}{Dependent variable: log real hourly labor income} \\
    \cmidrule(lr){2-5}
    & (1) & (2) & (3) & (4) \\
    \midrule
    Firm size: 2-3 & 0.184*** & 0.242*** & 0.209*** & 0.188*** \\
     & (0.034) & (0.031) & (0.027) & (0.029) \\
    Firm size: 4-5 & 0.363*** & 0.408*** & 0.361*** & 0.307*** \\
     & (0.034) & (0.030) & (0.029) & (0.038) \\
    Firm size: 6-10 & 0.499*** & 0.500*** & 0.422*** & 0.316*** \\
     & (0.029) & (0.033) & (0.037) & (0.050) \\
    Firm size: 11-19 & 0.625*** & 0.575*** & 0.457*** & 0.302*** \\
     & (0.034) & (0.040) & (0.042) & (0.050) \\
    Firm size: 20-30 & 0.687*** & 0.603*** & 0.462*** & 0.271*** \\
     & (0.046) & (0.045) & (0.046) & (0.051) \\
    Firm size: 31-50 & 0.759*** & 0.650*** & 0.487*** & 0.271*** \\
     & (0.056) & (0.049) & (0.047) & (0.048) \\
    Firm size: 51-100 & 0.818*** & 0.698*** & 0.518*** & 0.289*** \\
     & (0.063) & (0.051) & (0.048) & (0.047) \\
    Firm size: 101+ & 1.017*** & 0.817*** & 0.598*** & 0.359*** \\
     & (0.102) & (0.067) & (0.061) & (0.059) \\
    \midrule
    Gender, age, and education controls & No & No & Yes & Yes \\
    Formality control & No & No & No & Yes \\
    Sector-year fixed effects & No & Yes & Yes & Yes \\
    Observations & 5,268,514 & 5,268,514 & 5,268,514 & 5,268,514 \\
    $R^2$ & 0.213 & 0.271 & 0.376 & 0.386 \\
    \bottomrule
  \end{tabular}
  \vspace{0.3em}
  \begin{minipage}{0.95\textwidth}
  \footnotesize
  Notes: The omitted category is solo workers. All columns use the same estimation sample and GEIH expansion weights. Standard errors, clustered by sector, are reported in parentheses. Worker controls include a female-worker indicator, age, age squared, and education dummies. Column (4) adds labor formality and corresponds to Equation~(\ref{eq:common-premium}). All columns constrain size coefficients to be common across formality groups. Significance levels: * $p<0.10$, ** $p<0.05$, *** $p<0.01$.
  \end{minipage}
\end{table}

\begin{figure}[htbp]
  \centering
  \includegraphics[width=0.9\textwidth]{fig57_\figurelang.png}
  \caption{Common-premium benchmark relative to solo workers}
  \label{fig:firm-size-wage-premium}
\end{figure}
\begin{table}[htbp]
  \centering
  \caption{Firm-size wage premium regressions: 101+ workers as the reference category}
  \label{tab:firm-size-wage-premium-regressions-101-reference}
  \small
  \begin{tabular}{lcccc}
    \toprule
    & \multicolumn{4}{c}{Dependent variable: log real hourly labor income} \\
    \cmidrule(lr){2-5}
    & (1) & (2) & (3) & (4) \\
    \midrule
    Firm size: Solo & -1.028*** & -0.832*** & -0.613*** & -0.372*** \\
     & (0.102) & (0.068) & (0.062) & (0.060) \\
    Firm size: 2-3 & -0.838*** & -0.584*** & -0.398*** & -0.179*** \\
     & (0.095) & (0.056) & (0.046) & (0.042) \\
    Firm size: 4-5 & -0.656*** & -0.414*** & -0.244*** & -0.056* \\
     & (0.098) & (0.052) & (0.040) & (0.032) \\
    Firm size: 6-10 & -0.522*** & -0.323*** & -0.184*** & -0.049 \\
     & (0.092) & (0.051) & (0.037) & (0.032) \\
    Firm size: 11-19 & -0.394*** & -0.247*** & -0.147*** & -0.062 \\
     & (0.081) & (0.054) & (0.041) & (0.038) \\
    Firm size: 20-30 & -0.330*** & -0.217*** & -0.139*** & -0.090** \\
     & (0.069) & (0.050) & (0.038) & (0.038) \\
    Firm size: 31-50 & -0.258*** & -0.169*** & -0.113*** & -0.090** \\
     & (0.059) & (0.045) & (0.033) & (0.034) \\
    Firm size: 51-100 & -0.199*** & -0.121*** & -0.082*** & -0.072** \\
     & (0.051) & (0.036) & (0.027) & (0.028) \\
    \midrule
    Gender, age, and education controls & No & No & Yes & Yes \\
    Formality control & No & No & No & Yes \\
    Sector-year fixed effects & No & Yes & Yes & Yes \\
    Observations & 5,426,316 & 5,426,316 & 5,426,316 & 5,426,316 \\
    $R^2$ & 0.218 & 0.277 & 0.381 & 0.391 \\
    \bottomrule
  \end{tabular}
  \vspace{0.3em}
  \begin{minipage}{0.95\textwidth}
  \footnotesize
  Notes: Specifications as in Table~\ref{tab:firm-size-wage-premium-regressions}, with 101+ workers as the omitted category. Parentheses contain sector-clustered standard errors. * $p<0.10$, ** $p<0.05$, *** $p<0.01$.
  \end{minipage}
\end{table}

\begin{table}[htbp]
  \centering
  \caption{Firm-size wage premium regressions: 11--19 workers as the reference category}
  \label{tab:firm-size-wage-premium-regressions-11-19-reference}
  \small
  \begin{tabular}{lcccc}
    \toprule
    & \multicolumn{4}{c}{Dependent variable: log real hourly labor income} \\
    \cmidrule(lr){2-5}
    & (1) & (2) & (3) & (4) \\
    \midrule
    Firm size: Solo & -0.633*** & -0.585*** & -0.466*** & -0.310*** \\
     & (0.035) & (0.041) & (0.042) & (0.051) \\
    Firm size: 2-3 & -0.444*** & -0.337*** & -0.252*** & -0.117*** \\
     & (0.029) & (0.033) & (0.028) & (0.030) \\
    Firm size: 4-5 & -0.262*** & -0.167*** & -0.097*** & 0.006 \\
     & (0.029) & (0.021) & (0.021) & (0.020) \\
    Firm size: 6-10 & -0.128*** & -0.076*** & -0.037*** & 0.013 \\
     & (0.017) & (0.010) & (0.010) & (0.010) \\
    Firm size: 20-30 & 0.064*** & 0.030*** & 0.008 & -0.028*** \\
     & (0.015) & (0.009) & (0.006) & (0.005) \\
    Firm size: 31-50 & 0.136*** & 0.078*** & 0.034*** & -0.028*** \\
     & (0.026) & (0.015) & (0.010) & (0.008) \\
    Firm size: 51-100 & 0.196*** & 0.126*** & 0.065*** & -0.010 \\
     & (0.034) & (0.023) & (0.017) & (0.014) \\
    Firm size: 101+ & 0.394*** & 0.247*** & 0.147*** & 0.062 \\
     & (0.081) & (0.054) & (0.041) & (0.038) \\
    \midrule
    Gender, age, and education controls & No & No & Yes & Yes \\
    Formality control & No & No & No & Yes \\
    Sector-year fixed effects & No & Yes & Yes & Yes \\
    Observations & 5,426,316 & 5,426,316 & 5,426,316 & 5,426,316 \\
    $R^2$ & 0.218 & 0.277 & 0.381 & 0.391 \\
    \bottomrule
  \end{tabular}
  \vspace{0.3em}
  \begin{minipage}{0.95\textwidth}
  \footnotesize
  Notes: Specifications as in Table~\ref{tab:firm-size-wage-premium-regressions}, with 11--19 workers as the omitted category. Parentheses contain sector-clustered standard errors. * $p<0.10$, ** $p<0.05$, *** $p<0.01$.
  \end{minipage}
\end{table}


Tables~\ref{tab:firm-size-local-geography-robustness},
\ref{tab:firm-size-local-geography-robustness-101-reference}, and
\ref{tab:firm-size-local-geography-robustness-11-19-reference} add geographic
and occupational-position controls while retaining the common-premium
restriction. Their final columns give a 40.2 percent 101+/solo premium.
Unlike the preferred model, these regressions do not allow the formal and
informal size profiles to differ.

\begin{table}[htbp]
  \centering
  \caption{Common-premium estimates with geographic and occupational controls}
  \label{tab:firm-size-local-geography-robustness}
  \small
  \begin{tabular}{lcccc}
    \toprule
    & \multicolumn{4}{c}{Dependent variable: log real hourly labor income} \\
    \cmidrule(lr){2-5}
    & (1) & (2) & (3) & (4) \\
    \midrule
    Firm size: 2-3 & 0.188*** & 0.168*** & 0.167*** & 0.102*** \\
     & (0.011) & (0.011) & (0.011) & (0.013) \\
    Firm size: 4-5 & 0.307*** & 0.277*** & 0.273*** & 0.225*** \\
     & (0.014) & (0.014) & (0.014) & (0.015) \\
    Firm size: 6-10 & 0.316*** & 0.282*** & 0.278*** & 0.265*** \\
     & (0.018) & (0.017) & (0.017) & (0.016) \\
    Firm size: 11-19 & 0.302*** & 0.272*** & 0.267*** & 0.272*** \\
     & (0.018) & (0.017) & (0.017) & (0.015) \\
    Firm size: 20-30 & 0.271*** & 0.241*** & 0.237*** & 0.258*** \\
     & (0.019) & (0.017) & (0.017) & (0.015) \\
    Firm size: 31-50 & 0.271*** & 0.245*** & 0.240*** & 0.263*** \\
     & (0.018) & (0.017) & (0.017) & (0.015) \\
    Firm size: 51-100 & 0.289*** & 0.262*** & 0.257*** & 0.282*** \\
     & (0.018) & (0.017) & (0.017) & (0.015) \\
    Firm size: 101+ & 0.359*** & 0.338*** & 0.325*** & 0.338*** \\
     & (0.023) & (0.024) & (0.024) & (0.022) \\
    \midrule
    Worker and formality controls & Yes & Yes & Yes & Yes \\
    Occupational-position controls & No & No & No & Yes \\
    Fixed effects & $s\times t$ & $s\times t$, $d\times t$ & $s\times d\times t$ & $s\times d\times t$ \\
    Observations & 5,268,514 & 5,268,514 & 5,268,375 & 5,268,375 \\
    $R^2$ & 0.386 & 0.404 & 0.414 & 0.429 \\
    \bottomrule
  \end{tabular}
  \vspace{0.3em}
  \begin{minipage}{0.95\textwidth}
  \footnotesize
  Notes: The omitted category is solo workers. Column (4) corresponds to Equation~(\ref{eq:common-premium-geography}). $s$, $d$, and $t$ denote sector, department, and year. Parentheses contain sector-department-clustered standard errors. * $p<0.10$, ** $p<0.05$, *** $p<0.01$.
  \end{minipage}
\end{table}

\begin{table}[htbp]
  \centering
  \caption{Common-premium estimates with geographic and status-in-employment controls: 101+ workers as reference}
  \label{tab:firm-size-local-geography-robustness-101-reference}
  \small
  \begin{tabular}{lcccc}
    \toprule
    & \multicolumn{4}{c}{Dependent variable: log real hourly labor income} \\
    \cmidrule(lr){2-5}
    & (1) & (2) & (3) & (4) \\
    \midrule
    Firm size: Solo & -0.372*** & -0.352*** & -0.337*** & -0.343*** \\
     & (0.024) & (0.024) & (0.024) & (0.022) \\
    Firm size: 2-3 & -0.179*** & -0.178*** & -0.165*** & -0.242*** \\
     & (0.018) & (0.018) & (0.017) & (0.016) \\
    Firm size: 4-5 & -0.056*** & -0.066*** & -0.056*** & -0.117*** \\
     & (0.013) & (0.014) & (0.014) & (0.014) \\
    Firm size: 6-10 & -0.049*** & -0.063*** & -0.053*** & -0.078*** \\
     & (0.012) & (0.013) & (0.013) & (0.012) \\
    Firm size: 11-19 & -0.062*** & -0.072*** & -0.064*** & -0.071*** \\
     & (0.014) & (0.015) & (0.015) & (0.013) \\
    Firm size: 20-30 & -0.090*** & -0.100*** & -0.091*** & -0.082*** \\
     & (0.014) & (0.015) & (0.014) & (0.012) \\
    Firm size: 31-50 & -0.090*** & -0.096*** & -0.088*** & -0.076*** \\
     & (0.012) & (0.014) & (0.013) & (0.011) \\
    Firm size: 51-100 & -0.072*** & -0.079*** & -0.070*** & -0.058*** \\
     & (0.012) & (0.013) & (0.012) & (0.011) \\
    \midrule
    Worker and formality controls & Yes & Yes & Yes & Yes \\
    Status-in-employment controls & No & No & No & Yes \\
    Fixed effects & $s\times t$ & $s\times t$, $d\times t$ & $s\times d\times t$ & $s\times d\times t$ \\
    Observations & 5,426,316 & 5,426,316 & 5,426,174 & 5,426,174 \\
    $R^2$ & 0.391 & 0.410 & 0.420 & 0.435 \\
    \bottomrule
  \end{tabular}
  \vspace{0.3em}
  \begin{minipage}{0.95\textwidth}
  \footnotesize
  Notes: Specifications as in Table~\ref{tab:firm-size-local-geography-robustness}, with 101+ workers as the omitted category. $s$, $d$, and $t$ denote sector, department, and year. Parentheses contain sector-department-clustered standard errors. * $p<0.10$, ** $p<0.05$, *** $p<0.01$.
  \end{minipage}
\end{table}

\begin{table}[htbp]
  \centering
  \caption{Common-premium estimates with geographic and occupational controls: 11--19 workers as reference}
  \label{tab:firm-size-local-geography-robustness-11-19-reference}
  \small
  \begin{tabular}{lcccc}
    \toprule
    & \multicolumn{4}{c}{Dependent variable: log real hourly labor income} \\
    \cmidrule(lr){2-5}
    & (1) & (2) & (3) & (4) \\
    \midrule
    Firm size: Solo & -0.302*** & -0.272*** & -0.267*** & -0.272*** \\
     & (0.018) & (0.017) & (0.017) & (0.015) \\
    Firm size: 2-3 & -0.114*** & -0.104*** & -0.100*** & -0.169*** \\
     & (0.013) & (0.011) & (0.011) & (0.009) \\
    Firm size: 4-5 & 0.005 & 0.005 & 0.006 & -0.047*** \\
     & (0.008) & (0.008) & (0.008) & (0.007) \\
    Firm size: 6-10 & 0.014*** & 0.010** & 0.011** & -0.007 \\
     & (0.005) & (0.005) & (0.005) & (0.004) \\
    Firm size: 20-30 & -0.031*** & -0.031*** & -0.030*** & -0.014*** \\
     & (0.004) & (0.004) & (0.004) & (0.004) \\
    Firm size: 31-50 & -0.031*** & -0.027*** & -0.027*** & -0.008* \\
     & (0.005) & (0.005) & (0.005) & (0.004) \\
    Firm size: 51-100 & -0.014* & -0.010 & -0.010 & 0.010 \\
     & (0.007) & (0.007) & (0.007) & (0.007) \\
    Firm size: 101+ & 0.057*** & 0.066*** & 0.058*** & 0.066*** \\
     & (0.014) & (0.015) & (0.015) & (0.013) \\
    \midrule
    Worker and formality controls & Yes & Yes & Yes & Yes \\
    Occupational-position controls & No & No & No & Yes \\
    Fixed effects & $s\times t$ & $s\times t$, $d\times t$ & $s\times d\times t$ & $s\times d\times t$ \\
    Observations & 5,268,514 & 5,268,514 & 5,268,375 & 5,268,375 \\
    $R^2$ & 0.386 & 0.404 & 0.414 & 0.429 \\
    \bottomrule
  \end{tabular}
  \vspace{0.3em}
  \begin{minipage}{0.95\textwidth}
  \footnotesize
  Notes: Specifications as in Table~\ref{tab:firm-size-local-geography-robustness}, with 11--19 workers as the omitted category. $s$, $d$, and $t$ denote sector, department, and year. Parentheses contain sector-department-clustered standard errors. * $p<0.10$, ** $p<0.05$, *** $p<0.01$.
  \end{minipage}
\end{table}

\FloatBarrier

The simpler interacted comparison model is
\begin{equation}\label{eq:formality-interaction}
  \ln(w_{ist}) = \alpha + \widetilde{\phi}F_{ist}
  + \sum_{g\neq g_0}(\widetilde{\beta}_g+
       \widetilde{\theta}_g F_{ist})D_{ig}
  + X_{ist}'\gamma+\lambda_{s\times t}+\varepsilon_{ist}.
\end{equation}
Tildes distinguish these coefficients from the preferred model's coefficients.
This regression includes worker characteristics and sector-year effects, but
not occupational-position or department controls. Its inference clusters by
sector. The informal 101+/solo premium is 39.9 percent, and the formal
premium is 16.6 percent, with a sector-clustered \(p\)-value of 0.1001 for the
latter. Formal 101+ workers nevertheless earn 8.9--14.1 percent more than
workers in the smaller non-solo formal categories, with significant direct
contrasts. The associated formal-informal gaps and alternative normalizations
are reported below. These estimates are not the source of the preferred
income scenarios.

\begin{table}[htbp]
  \centering
  \caption{Formality-specific coefficients in the sector-year comparison model}
  \label{tab:formality-interaction-coefficients}
  \small
  \begin{tabular}{lccc}
    \toprule
    Firm size & $\hat{\beta}_g$ & $\hat{\theta}_g$ & $\hat{\beta}_g+\hat{\theta}_g$ \\
    \midrule
    2-3 & 0.197*** & -0.175** & 0.022 \\
     & (0.026) & (0.062) & (0.067) \\
    4-5 & 0.343*** & -0.318*** & 0.025 \\
     & (0.027) & (0.085) & (0.089) \\
    6-10 & 0.392*** & -0.369*** & 0.023 \\
     & (0.032) & (0.081) & (0.086) \\
    11-19 & 0.408*** & -0.373*** & 0.035 \\
     & (0.037) & (0.076) & (0.080) \\
    20-30 & 0.401*** & -0.378*** & 0.024 \\
     & (0.035) & (0.072) & (0.079) \\
    31-50 & 0.399*** & -0.358*** & 0.042 \\
     & (0.039) & (0.066) & (0.076) \\
    51-100 & 0.409*** & -0.340*** & 0.069 \\
     & (0.044) & (0.067) & (0.075) \\
    101+ & 0.336*** & -0.182** & 0.154 \\
     & (0.046) & (0.083) & (0.088) \\
    \bottomrule
  \end{tabular}
  \vspace{0.3em}
  \begin{minipage}{0.92\textwidth}
  \footnotesize
  Notes: Log-point coefficients from Equation~(\ref{eq:formality-interaction}), with solo workers as the reference within each formality group and tildes suppressed. The first and third columns are informal and formal size gaps. The middle column is the formal-minus-informal difference in those size gaps. Parentheses contain sector-clustered standard errors. * $p<0.10$, ** $p<0.05$, *** $p<0.01$.
  \end{minipage}
\end{table}

\begin{figure}[htbp]
  \centering
  \includegraphics[width=0.95\textwidth]{fig74_\figurelang.png}
  \caption{Firm-size premiums by formality in the sector-year comparison model}
  \label{fig:firm-size-formality-premium}
\end{figure}
\begin{table}[htbp]
  \centering
  \caption{Formal 101+ premiums in the sector-year comparison model}
  \label{tab:formal-101-contrasts}
  \small
  \begin{tabular}{lcc}
    \toprule
    Reference category & Premium & $p$-value \\
    \midrule
    Solo workers & 16.5\% & 0.082 \\
    2-3 & 14.2\% & $<0.001$ \\
    4-5 & 13.8\% & $<0.001$ \\
    6-10 & 13.9\% & $<0.001$ \\
    11-19 & 12.6\% & 0.004 \\
    20-30 & 13.9\% & 0.001 \\
    31-50 & 11.8\% & 0.002 \\
    51-100 & 8.9\% & 0.004 \\
    \bottomrule
  \end{tabular}
  \vspace{0.3em}
  \begin{minipage}{0.92\textwidth}
  \footnotesize
  Notes: Each row compares formal workers in firms with 101 or more workers with formal workers in the reference firm-size category. Estimates come from the comparison specification in Equation~(\ref{eq:formality-interaction}) with gender, age, age squared, education, and sector-year fixed effects. Premiums are computed as $100\times[\exp(\hat{\beta})-1]$ from the relevant linear contrast.
  \end{minipage}
\end{table}

\begin{table}[htbp]
  \centering
  \caption{Formal-informal wage gaps in the sector-year comparison model}
  \label{tab:formal-wage-gap-by-size}
  \small
  \begin{tabular}{lccc}
    \toprule
    Firm size & Log gap & Premium & $p$-value \\
    \midrule
    Solo & 0.522*** & 68.6\% & $<0.001$ \\
     & (0.063) &  &  \\
    2-3 & 0.347*** & 41.5\% & $<0.001$ \\
     & (0.018) &  &  \\
    4-5 & 0.204*** & 22.6\% & $<0.001$ \\
     & (0.035) &  &  \\
    6-10 & 0.153*** & 16.6\% & $<0.001$ \\
     & (0.031) &  &  \\
    11-19 & 0.149*** & 16.1\% & $<0.001$ \\
     & (0.024) &  &  \\
    20-30 & 0.144*** & 15.5\% & $<0.001$ \\
     & (0.018) &  &  \\
    31-50 & 0.164*** & 17.9\% & $<0.001$ \\
     & (0.014) &  &  \\
    51-100 & 0.182*** & 20.0\% & $<0.001$ \\
     & (0.019) &  &  \\
    101+ & 0.340*** & 40.5\% & $<0.001$ \\
     & (0.043) &  &  \\
    \bottomrule
  \end{tabular}
  \vspace{0.3em}
  \begin{minipage}{0.92\textwidth}
  \footnotesize
  Notes: Formal-minus-informal contrasts from Equation~(\ref{eq:formality-interaction}) at each firm size. Percentage premiums use informal workers as the denominator. Parentheses contain sector-clustered standard errors. * $p<0.10$, ** $p<0.05$, *** $p<0.01$.
  \end{minipage}
\end{table}

\begin{table}[htbp]
  \centering
  \caption{Firm-size coefficients by formality status: 101+ workers as reference}
  \label{tab:formality-interaction-coefficients-101-reference}
  \small
  \begin{tabular}{lccc}
    \toprule
    Firm size & $\hat{\beta}_g$ & $\hat{\theta}_g$ & $\hat{\beta}_g+\hat{\theta}_g$ \\
    \midrule
    Solo & -0.336*** & 0.182** & -0.154 \\
     & (0.046) & (0.083) & (0.088) \\
    2-3 & -0.139*** & 0.007 & -0.132*** \\
     & (0.034) & (0.047) & (0.040) \\
    4-5 & 0.007 & -0.136** & -0.129*** \\
     & (0.031) & (0.048) & (0.033) \\
    6-10 & 0.056* & -0.187*** & -0.131*** \\
     & (0.029) & (0.046) & (0.035) \\
    11-19 & 0.072** & -0.191*** & -0.119** \\
     & (0.027) & (0.044) & (0.041) \\
    20-30 & 0.066** & -0.196*** & -0.130*** \\
     & (0.024) & (0.042) & (0.040) \\
    31-50 & 0.064** & -0.175*** & -0.112*** \\
     & (0.022) & (0.041) & (0.036) \\
    51-100 & 0.073*** & -0.158*** & -0.085** \\
     & (0.021) & (0.039) & (0.030) \\
    \bottomrule
  \end{tabular}
  \vspace{0.3em}
  \begin{minipage}{0.92\textwidth}
  \footnotesize
  Notes: Log-point coefficients from Equation~(\ref{eq:formality-interaction}), with 101+ workers as the reference within each formality group and tildes suppressed. The first and third columns are informal and formal size gaps. The middle column is the formal-minus-informal difference in those size gaps. Parentheses contain sector-clustered standard errors. * $p<0.10$, ** $p<0.05$, *** $p<0.01$.
  \end{minipage}
\end{table}

\begin{table}[htbp]
  \centering
  \caption{Firm-size coefficients by formality status: 11--19 workers as reference}
  \label{tab:formality-interaction-coefficients-11-19-reference}
  \small
  \begin{tabular}{lccc}
    \toprule
    Firm size & $\hat{\beta}_g$ & $\hat{\theta}_g$ & $\hat{\beta}_g+\hat{\theta}_g$ \\
    \midrule
    Solo & -0.415*** & 0.372*** & -0.043 \\
     & (0.038) & (0.075) & (0.079) \\
    2-3 & -0.212*** & 0.191*** & -0.021 \\
     & (0.025) & (0.022) & (0.024) \\
    4-5 & -0.063*** & 0.049** & -0.013 \\
     & (0.019) & (0.018) & (0.018) \\
    6-10 & -0.016 & 0.002 & -0.014 \\
     & (0.010) & (0.012) & (0.012) \\
    20-30 & -0.005 & -0.004 & -0.009* \\
     & (0.007) & (0.009) & (0.005) \\
    31-50 & -0.009 & 0.019 & 0.010 \\
     & (0.009) & (0.012) & (0.008) \\
    51-100 & 0.004 & 0.032 & 0.036** \\
     & (0.023) & (0.026) & (0.015) \\
    101+ & -0.074** & 0.198*** & 0.124*** \\
     & (0.027) & (0.044) & (0.040) \\
    \bottomrule
  \end{tabular}
  \vspace{0.3em}
  \begin{minipage}{0.92\textwidth}
  \footnotesize
  Notes: Log-point coefficients from Equation~(\ref{eq:formality-interaction}), with 11--19 workers as the reference within each formality group and tildes suppressed. The first and third columns are informal and formal size gaps. The middle column is the formal-minus-informal difference in those size gaps. Parentheses contain sector-clustered standard errors. * $p<0.10$, ** $p<0.05$, *** $p<0.01$.
  \end{minipage}
\end{table}

\FloatBarrier

The earlier dynamic benchmark interacts size with year but constrains the
premium to be common across formality groups:
\begin{equation}\label{eq:dynamic-benchmark}
  \ln(w_{ist}) = \alpha
  + \sum_{\tau}\sum_{g\neq g_0} b_{g\tau}D_{ig}\mathbf{1}\{t=\tau\}
  + X_{ist}'\gamma+\phi F_{ist}
  + \lambda_{s\times t}+\varepsilon_{ist}.
\end{equation}
Its linear-trend counterpart is
\begin{equation}\label{eq:trend-benchmark}
  \ln(w_{ist}) = \alpha
  + \sum_{g\neq g_0}(b_g+\delta_g Trend_t)D_{ig}
  + X_{ist}'\gamma+\phi F_{ist}
  + \lambda_{s\times t}+\varepsilon_{ist},
\end{equation}
where \(Trend_t=t-2008\). Two-sided tests of \(H_0:\delta_g=0\) do not reject
zero trends at 5 percent. These common-premium trends differ from the preferred
within-formality endpoint contrasts. The alternative-reference trend tables
use 101+ and 11--19 workers as the respective bases.

\begin{figure}[htbp]
  \centering
  \includegraphics[width=0.95\textwidth]{fig59_\figurelang.png}
  \caption{Common-premium benchmark over time}
  \label{fig:firm-size-wage-premium-over-time}
\end{figure}
\begin{table}[htbp]
  \centering
  \caption{Testing whether the firm-size wage premium changed over time}
  \label{tab:firm-size-premium-trend-test}
  \small
  \begin{tabular}{lcccc}
    \toprule
    Firm size & Annual trend & S.E. & $p$-value & Implied 2008--2025 change \\
    \midrule
    2-3 & 0.0026 & (0.0026) & 0.337 & 4.4\% \\
    4-5 & 0.0030 & (0.0029) & 0.314 & 5.2\% \\
    6-10 & 0.0022 & (0.0031) & 0.498 & 3.8\% \\
    11-19 & 0.0026 & (0.0037) & 0.485 & 4.6\% \\
    20-30 & 0.0014 & (0.0040) & 0.733 & 2.4\% \\
    31-50 & 0.0024 & (0.0041) & 0.556 & 4.2\% \\
    51-100 & 0.0025 & (0.0041) & 0.547 & 4.4\% \\
    101+ & 0.0038 & (0.0041) & 0.372 & 6.7\% \\
    \bottomrule
  \end{tabular}
  \vspace{0.3em}
  \begin{minipage}{0.95\textwidth}
  \footnotesize
  Notes: Annual changes in log-wage gaps from Equation~(\ref{eq:trend-benchmark}), expressed relative to solo workers. The final column is the compounded 2008--2025 percentage change in that wage ratio, not a percentage-point change in the premium. Two-sided tests of zero trends use sector-clustered standard errors.
  \end{minipage}
\end{table}

\begin{table}[htbp]
  \centering
  \caption{Testing whether wage gaps relative to 101+ workers changed over time}
  \label{tab:firm-size-premium-trend-test-101-reference}
  \small
  \begin{tabular}{lcccc}
    \toprule
    Firm size & Annual trend & S.E. & $p$-value & Implied 2008--2025 change \\
    \midrule
    Solo & -0.0045 & (0.0043) & 0.313 & -7.3\% \\
    2-3 & -0.0012 & (0.0023) & 0.612 & -2.0\% \\
    4-5 & -0.0013 & (0.0023) & 0.587 & -2.1\% \\
    6-10 & -0.0023 & (0.0017) & 0.190 & -3.9\% \\
    11-19 & -0.0017 & (0.0012) & 0.175 & -2.9\% \\
    20-30 & -0.0026 & (0.0009) & 0.012 & -4.3\% \\
    31-50 & -0.0017 & (0.0009) & 0.065 & -2.9\% \\
    51-100 & -0.0018 & (0.0013) & 0.181 & -3.1\% \\
    \bottomrule
  \end{tabular}
  \vspace{0.3em}
  \begin{minipage}{0.95\textwidth}
  \footnotesize
  Notes: Annual changes in log-wage gaps from Equation~(\ref{eq:trend-benchmark}), expressed relative to 101+ workers. The final column is the compounded 2008--2025 percentage change in that wage ratio, not a percentage-point change in the premium. Two-sided tests of zero trends use sector-clustered standard errors.
  \end{minipage}
\end{table}

\begin{table}[htbp]
  \centering
  \caption{Testing whether wage gaps relative to 11--19 workers changed over time}
  \label{tab:firm-size-premium-trend-test-11-19-reference}
  \small
  \begin{tabular}{lcccc}
    \toprule
    Firm size & Annual trend & S.E. & $p$-value & Implied 2008--2025 change \\
    \midrule
    Solo & -0.0027 & (0.0039) & 0.491 & -4.5\% \\
    2-3 & 0.0005 & (0.0016) & 0.746 & 0.9\% \\
    4-5 & 0.0005 & (0.0014) & 0.736 & 0.8\% \\
    6-10 & -0.0006 & (0.0011) & 0.611 & -1.0\% \\
    20-30 & -0.0009 & (0.0010) & 0.392 & -1.5\% \\
    31-50 & 0.0000 & (0.0012) & 0.988 & 0.0\% \\
    51-100 & -0.0001 & (0.0014) & 0.940 & -0.2\% \\
    101+ & 0.0017 & (0.0012) & 0.175 & 3.0\% \\
    \bottomrule
  \end{tabular}
  \vspace{0.3em}
  \begin{minipage}{0.95\textwidth}
  \footnotesize
  Notes: Annual changes in log-wage gaps from Equation~(\ref{eq:trend-benchmark}), expressed relative to 11--19 workers. The final column is the compounded 2008--2025 percentage change in that wage ratio, not a percentage-point change in the premium. Two-sided tests of zero trends use sector-clustered standard errors.
  \end{minipage}
\end{table}


The next two figures estimate Equation~(\ref{eq:formality-interaction})
separately by year, replacing sector-year effects with sector effects.
They omit occupational-position and department controls and use
sector-clustered inference. The 2025 formal 101+/solo estimate is
10.5 percent and its interval includes zero, whereas the informal estimate
is 53.6 percent. These simpler profiles are retained to make the specification
comparison reproducible.

\begin{figure}[htbp]
  \centering
  \includegraphics[width=0.95\textwidth]{fig76_\figurelang.png}
  \caption{Sector-control comparison of formality-specific premiums, 2008 and 2025}
  \label{fig:adjusted-formality-endpoints}
\end{figure}

\begin{figure}[htbp]
  \centering
  \includegraphics[width=0.95\textwidth]{fig75_\figurelang.png}
  \caption{Sector-control comparison of large-firm premiums by formality over time}
  \label{fig:large-firm-formality-premium-over-time}
\end{figure}
\FloatBarrier

The earlier log-log comparison uses the same assigned sizes as the preferred
elasticity exercise but retains the simpler controls:
\begin{equation}\label{eq:elasticity}
  \ln(w_{ist}) = \alpha+\eta\ln(\widetilde{Size}_{ist})
  + X_{ist}'\gamma+\phi F_{ist}
  + \lambda_{s\times t}+\varepsilon_{ist}.
\end{equation}
The raw model includes assigned log size only. The controlled model includes
worker characteristics, formality, and sector-year effects. Formal and
informal slopes in this earlier exercise come from separate regressions,
which omit the constant formality indicator. All use sector clustering.

\begin{table}[htbp]
  \centering
  \caption{Benchmarking Colombia's approximate firm-size wage elasticity}
  \label{tab:firm-size-elasticity-benchmark}
  \small
  \begin{tabular}{p{0.30\textwidth}p{0.40\textwidth}cc}
    \toprule
    Source & Specification & Elasticity & S.E. \\
    \midrule
    Diegmann et al. (2026) & No worker-heterogeneity controls & 0.062 & n.a. \\
    Diegmann et al. (2026) & Observable worker controls & 0.036 & n.a. \\
    \midrule
    Colombia (GEIH) & No controls & 0.203*** & (0.019) \\
    Colombia (GEIH) & Full sample & 0.071*** & (0.013) \\
    Colombia (GEIH) & Formal workers & 0.026** & (0.011) \\
    Colombia (GEIH) & Informal workers & 0.130*** & (0.015) \\
    \bottomrule
  \end{tabular}
  \vspace{0.3em}
  \begin{minipage}{0.95\textwidth}
  \footnotesize
  Notes: Literature means are from Table 1 of \citet{DiegmannMullerSchoefer2026}. The controlled full-sample row uses Equation~(\ref{eq:elasticity}). Formal and informal rows re-estimate it on separate samples. Standard errors are clustered by sector. * $p<0.10$, ** $p<0.05$, *** $p<0.01$.
  \end{minipage}
\end{table}


\FloatBarrier
We also ask whether two parameters per formality group can summarize the
eight size contrasts in the preferred categorical model.
\citet{FairDominguez1991} restrict coefficients on age-group population
shares to a quadratic polynomial and normalize them to sum to zero.
We adapt their polynomial restriction to firm-size coefficients, using log
assigned size rather than a category index. We retain solo work as the
reference instead of imposing a sum-zero normalization. This changes the
functional form, not the identification of the wage relationship.

Let \(z_i=\ln(\widetilde{Size}_i)\), using the assigned sizes in
Section~\ref{sec:methodology}. Write \(a_1\) and \(a_2\) for the informal
coefficients on log assigned size and its square. The corresponding
formal-minus-informal differences are \(b_1\) and \(b_2\), with \(k\)
indexing these two powers. The quadratic model is
\begin{equation}\label{eq:polynomial-size-profile}
  \ln(w_{iosdt})=\alpha+\phi F_i
  +\sum_{k=1}^{2}(a_k+b_k F_i)z_i^k
  +X_i'\gamma+\rho_o+\lambda_{s\times d\times t}
  +\varepsilon_{iosdt}.
\end{equation}
The sample, worker and occupational-position controls, fixed effects, weights,
and baseline clustering are the same as in Equation~(\ref{eq:main}). For status
\(f\in\{0,1\}\), write \(a_{kf}=a_k+b_k f\). The log-wage contrast
relative to solo work is \(q_f(z)=a_{1f}z+a_{2f}z^2\), which is zero at
the assigned solo size of one. Its implied elasticity is
\(a_{1f}+2a_{2f}z\), not a constant elasticity. The curve between assigned
sizes is imposed by the quadratic rather than observed in the data.

\begin{table}[htbp]
  \centering
  \caption{Quadratic size profiles with preferred controls}
  \label{tab:polynomial-size-profile}
  \small
  \begin{tabular}{lcc}
    \toprule
    Term & Formal workers & Informal workers \\[0.25em]
    \midrule
    Log assigned size ($\widehat{a}_{1f}$) & \shortstack{-0.0086\\{\scriptsize (0.0147)}} & \shortstack{0.1956***\\{\scriptsize (0.0134)}} \\[0.25em]
    Squared log assigned size ($\widehat{a}_{2f}$) & \shortstack{0.0089***\\{\scriptsize (0.0021)}} & \shortstack{-0.0251***\\{\scriptsize (0.0026)}} \\[0.25em]
    \midrule
    Implied 101+/solo premium (\%) & 19.8 & 41.9 \\[0.25em]
    \bottomrule
  \end{tabular}
  \vspace{0.3em}
  \begin{minipage}{0.95\textwidth}
  \footnotesize
  Notes: Both columns come from Equation~(\ref{eq:polynomial-size-profile}) on the preferred sample, with 101+ assigned 150 workers. The last row reports percentage premiums relative to solo work within each formality group. Parentheses contain sector-department-clustered standard errors. * $p<0.10$, ** $p<0.05$, *** $p<0.01$.
  \end{minipage}
\end{table}


Table~\ref{tab:polynomial-size-profile} reports the quadratic coefficients
when 101+ is assigned 150 workers. Figure~\ref{fig:polynomial-size-profile}
overlays the resulting curves and pointwise 95 percent confidence bands on
the preferred categorical estimates and their intervals. The informal curve
rises rapidly at small sizes and then bends downward, while the formal curve
has a much smaller initial decline followed by an increase. The quadratic
does not reproduce every categorical comparison. For formal 2--3 workers,
it predicts a gap of approximately zero relative to solo work, compared with -7.2 percent
in the categorical model. Its informal 101+/solo premium is 41.9 percent,
compared with 37.0 percent in the categorical model. The corresponding formal
101+/solo estimates are 19.8 and 17.2 percent.

\begin{figure}[htbp]
  \centering
  \includegraphics[width=0.95\textwidth]{preferred_quadratic_\figurelang.png}
  \caption{Quadratic and categorical firm-size profiles with preferred controls}
  \label{fig:polynomial-size-profile}
\end{figure}

Table~\ref{tab:quadratic-preferred-sensitivity} tests the functional-form
restriction and varies the size assigned to 101+. A quadratic constrains
eight within-formality categorical contrasts to depend on two coefficients,
imposing six restrictions per group and twelve jointly. We test those linear
restrictions using the unrestricted categorical estimates and their full
clustered covariance matrix. With 150 assigned to 101+, the sector-department
Wald \(F\)-statistics are 23.27 for formal workers and 22.34 for informal
workers, with six numerator and 558 denominator degrees of freedom.
Both tests reject the quadratic restriction at 1 percent, as does the joint
test.

\begin{table}[htbp]
  \centering
  \caption{Top-bin sensitivity and tests of the quadratic restriction}
  \label{tab:quadratic-preferred-sensitivity}
  \small
  \begin{tabular}{rrrrrr}
    \toprule
    Size for 101+ & Formal (\%) & Informal (\%) & $p_F$ & $p_I$ & $p_{Joint}$ \\[0.25em]
    \midrule
    \multicolumn{6}{l}{\textit{Sector-department clustering}} \\
    101 & 17.9 & 43.2 & $<0.001$ & $<0.001$ & $<0.001$ \\[0.25em]
    150 & 18.9 & 41.9 & $<0.001$ & $<0.001$ & $<0.001$ \\[0.25em]
    250 & 19.8 & 40.5 & $<0.001$ & $<0.001$ & $<0.001$ \\[0.25em]
    \midrule
    \multicolumn{6}{l}{\textit{Sector clustering}} \\
    101 & 17.9 & 43.2 & $<0.001$ & $<0.001$ & $<0.001$ \\[0.25em]
    150 & 18.9 & 41.9 & $<0.001$ & $<0.001$ & $<0.001$ \\[0.25em]
    250 & 19.8 & 40.5 & $<0.001$ & 0.001 & $<0.001$ \\[0.25em]
    \bottomrule
  \end{tabular}
  \vspace{0.3em}
  \begin{minipage}{0.95\textwidth}
  \footnotesize
  Notes: Formal and informal percentages are the implied 101+/solo premiums from the quadratic model, re-estimated for each assigned top-bin size. All other size assignments, observations, controls, and weights remain fixed. The $p$-values test whether the unrestricted categorical profiles from Equation~(\ref{eq:main}) lie on the specified quadratic. There are six restrictions per formality group and twelve jointly. Cluster-robust Wald statistics divided by the number of restrictions use an $F_{q,G-1}$ reference distribution. There are 559 sector-department clusters or 17 sector clusters. These are approximate tests, especially with only 17 clusters. The repeated percentages emphasize that changing clustering does not change point estimates. The value 150 is the baseline assignment, not an observed mean for 101+.
  \end{minipage}
\end{table}


Assigning 101, 150, or 250 workers to the open-ended category gives formal
101+/solo premiums of 18.7, 19.8, and 20.8 percent, respectively.
The informal counterparts are 43.2, 41.9, and 40.5 percent. The quadratic
restrictions are rejected at 1 percent under every assignment with either
sector-department or sector clustering. The latter tests use only 17 clusters
and should be read as approximate sensitivity checks. The shallow formal dip
under the baseline assignment lies between assigned sizes one and 2.5.
It disappears when 101+ is assigned 250, so it is not stable evidence of a
U-shaped conditional relationship. Confidence bands hold the size assignments
fixed and do not incorporate this uncertainty. We retain the categorical model as
the preferred specification.
\FloatBarrier

